% !TeX TS-program = xelatex
% Standalone, one-page A4 resume. Compile with XeLaTeX or Tectonic.
\documentclass[10pt,a4paper]{article}
\usepackage[margin=12mm]{geometry}
\usepackage[T1]{fontenc}
\usepackage{tgtermes}
\usepackage{enumitem}
\usepackage{titlesec}
\usepackage{xcolor}
\usepackage[colorlinks=true,urlcolor=linkblue]{hyperref}

\definecolor{linkblue}{HTML}{225A85}
\hypersetup{pdftitle={Jiapeng Tang - AI Infrastructure Resume},pdfauthor={Jiapeng Tang}}
\pagestyle{empty}
\setlength{\parindent}{0pt}
\setlength{\parskip}{0pt}
\setlength{\emergencystretch}{1em}
\titleformat{\section}{\large\bfseries}{}{0pt}{}[\titlerule]
\titlespacing*{\section}{0pt}{6pt}{4pt}
\setlist[itemize]{leftmargin=12pt,label=\textbullet,topsep=2pt,itemsep=1.5pt,parsep=0pt,partopsep=0pt}

\newcommand{\entry}[2]{\par\noindent\textbf{#1}\hfill\mbox{#2}\par}
\newcommand{\pr}[2]{\href{https://github.com/#1/pull/#2}{\##2}}

\begin{document}
\fontsize{10.5pt}{12pt}\selectfont
\begin{center}
  {\LARGE\bfseries Jiapeng Tang}\\[3pt]
  AI Infrastructure \textbar{} LLM Inference \& GPU Kernel Optimization\\[2pt]
  +86 158-0628-1370 \textbar{}
  \href{mailto:3081129260@qq.com}{3081129260@qq.com} \textbar{}
  \href{https://github.com/tangzzycc}{github.com/tangzzycc}
\end{center}
\vspace{-7pt}

\section{Education}
\entry{Xi'an Jiaotong University}{Sep 2024 -- Jun 2027 (expected)}
\textit{Master's Degree in Software Engineering}\par
\textbf{GPA:} 3.82/4.0; multiple academic scholarships; \textbf{English:} CET-6.\par
\textbf{Research:} On-device LLM deployment,
mixed-precision inference, and NPU kernel optimization on Qualcomm Snapdragon.
\par\vspace{3pt}
\entry{China University of Petroleum (East China)}{Sep 2020 -- Jun 2024}
\textit{Bachelor's Degree in Software Engineering}\par

\section{Experience}
\entry{AMD\hspace{1em}Physical AI Intern}{Jul 2026 -- Present}
\textbf{vLLM ROCm Model Integration and Optimization}
\begin{itemize}
  \item Integrated CogAgent-18B multimodal inference into vLLM V1 on AMD RDNA 3.5 GPUs.
  Adapted dual vision encoders and per-layer visual cross-attention, enabled \texttt{torch.compile} and HIP Graph,
  and aligned BF16 outputs with Hugging Face.
  \item Enabled AWQ W4A16 inference and profiled the execution path with ROCm Profiler.
  Compared with the initial Hugging Face BF16 eager baseline, reduced TTFT by \textbf{47.4\%}, achieved
  \textbf{5.28$\times$ Decode throughput}, and reduced end-to-end latency by \textbf{76.1\%}.
  \item Implemented and tuned Triton LoRA Shrink/Expand kernels for small-batch multi-LoRA inference,
  increasing Decode throughput by \textbf{18.6--24.8\%} and reducing end-to-end latency by up to \textbf{19.0\%}.
  \item Optimized ROCm/RDNA GEMM and W4A16 kernel dispatch and fixed sliding-window Attention boundaries.
  Merged vLLM PRs: \pr{vllm-project/vllm}{53283}, \pr{vllm-project/vllm}{52619}, \pr{vllm-project/vllm}{59550}.
\end{itemize}
\textbf{FastFlowLM LoRA Inference Support}
\begin{itemize}
  \item Extended FastFlowLM's NPU inference framework with LoRA support for Gemma4-E models.
  Implemented PEFT adapter weight conversion, hardware-friendly layout transformations, and metadata generation
  to prepare adapters for NPU execution.
  \item Implemented runtime adapter loading and adapted Prefill/Decode AIE kernels for LoRA computation.
  Decoupled adapter loading and computation from base Q4NX weights and supported multiple common LoRA ranks.
\end{itemize}

\textbf{On-device Inference Framework Research}
\begin{itemize}
  \item Studied the architecture, inference workflows, and deployment pipelines of GenieX (Qualcomm QAIRT/QNN)
  and NVIDIA TensorRT-Edge-LLM for on-device LLM inference.
\end{itemize}

\section{Projects}
\entry{vLLM Adaptive KV Offload\hspace{1em}Independent Developer}{Jun 2026 -- Jul 2026}
\begin{itemize}
  \item Addressed high TTFT for short and medium requests that hit external KV cache but still wait in storage
  queues. Implemented a Load-vs-Recompute policy in vLLM V1 that compares estimated KV load and re-Prefill
  completion costs and dynamically selects cache reuse or recomputation per request.
  \item Used GPU profiling and I/O traces to analyze KV promotion, H2D transfers, and GPU Prefill timelines,
  identifying storage queuing and compute contention. Improved concurrent load/recompute scheduling and added
  recompute admission control and queued-I/O cancellation to reduce redundant transfers and extra Prefill work.
  \item Benchmarked Qwen3.6-35B-A3B-FP8 on 4$\times$RTX 5090 (TP4) under a storage-bound mixed workload.
  For policy-selected 4K latency-sensitive requests, reduced TTFT P50/P99 by \textbf{54.7\%/36.4\%};
  overall throughput slightly exceeded the baseline.
\end{itemize}
\entry{Qualcomm Snapdragon LLM Optimization\hspace{1em}Research Group Project}{Jan 2026 -- Apr 2026}
\begin{itemize}
  \item Deployed Qwen3-1.7B GPTQ W4A16 on Snapdragon 8 Gen 3 and the SA8797 automotive platform through
  ExecuTorch/QNN. Adapted quantization and graph compilation to import offline GPTQ weights and replace
  Linear layers, bypassing default online PTQ calibration to avoid requantizing the weights.
  \item Implemented a custom HVX lookup-table Linear operator, integrating GPTQ weight unpacking, tiled
  layout transformations, and activation lookup-table computation into the QNN backend. Combined projection
  fusion, MHA2SHA graph optimization, and 8-bit KV cache to reduce scheduling and memory overhead,
  improving Decode throughput by approximately \textbf{15\%} on Snapdragon 8 Gen 3.
  \item Extended the ExecuTorch/QNN runtime for Qwen3.5 hybrid Attention. Redesigned model state interfaces
  and Prefill/Decode state transfer to unify KV cache and linear-Attention state management; adapted model export,
  QNN graph compilation, and the device runner to enable end-to-end Qwen3.5-0.8B FP16 inference.
\end{itemize}

\section{Technical Skills}
\begin{itemize}
  \item \textbf{Programming \& Engineering:} C/C++, Python; Linux development/debugging; CMake, Docker, Git, Codex.
  \item \textbf{Inference \& Deployment:} vLLM, ExecuTorch/QNN; Continuous Batching, PagedAttention,
  quantization, speculative decoding; on-device deployment on Qualcomm and AMD platforms.
  \item \textbf{Kernels \& Profiling:} CUDA, HIP, Triton; GPU architecture and memory, kernel fusion and custom
  operators; performance analysis with Nsight Compute/Systems and ROCm Profiler.
\end{itemize}
\end{document}
